\section{Discussion}\label{sec:discussion}

% Every statement about results refers to Section 5 and uses its generated numbers. Statements
% about native models are limited to properties of the model family; what the runs of our
% configurations show beyond that is in Section S3.

\paragraph{What a patch model is for.} The answer of a scenario depends on the model of the
patch exactly where it depends on how the pressure is distributed: the friction torque of a
spinning face, the tilt of an eccentric load, the angle at which a body tips over a soft
support, the force that holds a press fit. A point contact represents a face by a few points
with a stiffness in N/m; it gives a disk the friction torque of a ring unless the user supplies
a torsional coefficient per geometry, and it has no areal stiffness against which class C could
be asked. Complementarity contact has no compliance at all, so the same questions do not apply
to it, and it is the natural choice where rigidity is the right model: a body sliding or
rolling on one plane, a body held by friction, stiff contacts on light bodies. On those
problems a patch model adds nothing, and the two laws pay for their foundation with the rolling
resistance of its damping, with its compliance and with the step bound of their explicit
coupling. Table~\ref{tab:decision} sets the situations against the model property that each one
needs.

\begin{table}[t]
\caption{Which property of a contact model each situation needs, as far as classes A and B of
the battery show}
\label{tab:decision}
\footnotesize
\begin{tabular}{@{}p{0.30\linewidth}p{0.33\linewidth}p{0.31\linewidth}@{}}
\toprule
Situation & What it needs & Evidence \\
\midrule
A rigid body sliding or rolling on one plane; sphere and cylinder contacts &
  Nothing beyond point or complementarity contact &
  Both laws within 1\,\% in \result{within/oneplane/distributed+gonthier/1} such scenarios
  (Table~\ref{tab:errors-A}); no gain from the patch \\
A flat face that spins, or slides and spins &
  A pressure distributed over the face: \law{distributed} or a pressure field; \law{gonthier}
  where its lumping error is acceptable &
  \scn{disk}: \result{err/disk/distributed}; \law{gonthier} \result{err/disk/gonthier};
  the rim torque of point contact (Section~\ref{sec:classB}) \\
Curved conforming contact (a journal, a conical seat) &
  Local normals: a per-face foundation or a pressure field. A finer mesh is not a remedy for
  \law{distributed} &
  \scn{journal}: \law{distributed} \result{err/journal/distributed/mesh32},
  \result{err/journal/distributed}, \result{err/journal/distributed/mesh128} at 32, 64 and 128
  segments; \law{pcm} \result{err/journal/pcm/mesh128} at 128 \\
Bodies or grounds cut into convex pieces &
  A wrench independent of the cut: \law{distributed} &
  \scn{pieces64}: \result{err/pieces64/distributed}; \law{gonthier}
  \result{err/pieces64/gonthier}, \law{pcm} \result{err/pieces64/pcm} \\
A body held by friction near its limit &
  A stick state, or a complementarity formulation &
  \scn{push}, \scn{wedge}: the laws below 1\,\%; \law{pcm} \result{err/push/pcm} and
  \result{err/wedge/pcm} \\
Stiff contacts on light bodies, large mass ratios &
  A formulation implicit in the contact force; not the explicit laws beyond
  Eq.~\eqref{eq:stability} &
  \scn{massratio}: the laws lost at 0.5\,ms and above, held from 0.2\,ms \\
Long rolling &
  No damping torque in the patch: rigid contact, or the laws with little damping or a stiff
  foundation &
  \scn{turntable}: \result{err/turntable/distributed};
  \result{err/turntable/distributed/stiff} with $K\times10$ \\
Only surface meshes are available &
  \law{gonthier} (any closed mesh) or \law{distributed} (convex pieces) &
  No volume mesh or pressure field is needed \\
\botrule
\end{tabular}
\end{table}

\paragraph{Lumped, distributed and per-face foundations.} The three are the same normal law.
\law{gonthier} lumps it, which is exact for a flat patch in the normal wrench and costs
\result{err/disk/gonthier} to \result{err/contensou/gonthier} where a face spins. \law{distributed} integrates it, which removes
that error and makes the result independent of the triangulation and of the cut, at a higher cost (medians of
\result{cost/gonthier/median} and \result{cost/distributed/median} of real time). \law{pcm} samples it once per element: what separates it from
\law{distributed} is the friction on coarse flat faces, which a finer mesh would repair, and
the lack of a stick state, which it would not. On curved conforming contact \law{pcm} is the
better of the two, because its depth follows the local normal: exact integration on each face
with a depth measured locally is the pressure field of hydroelastic contact reached from the
other side.

\paragraph{Against pressure-field contact.} With $E = K H$, hydroelastic contact is the same
pressure model as \law{distributed} near the surface, and the comparison is one of
discretizations. \law{distributed} takes surface meshes, needs convex pieces, clips the overlap
exactly and integrates pressure and friction on it in closed form, so that one triangle per
flat face is enough; its contact has one normal per pair of pieces and its coupling is
explicit. Hydroelastic contact needs a pressure field in every compliant body and has no
contact between two rigid geometries; its contact surface is curved, with no representative
normal, so it converges on curved conforming contact where \law{distributed} does not, and its
solver is implicit, so it has no step bound of the kind of Eq.~\eqref{eq:stability}; it pays
for the contact surface between tetrahedral fields and for one constraint per polygon. A user
who needs non-convex geometry, stiff contact or large steps is better served by it; a user who
has surface meshes, moderate stiffness and flat patches that spin, on convex or decomposed
bodies, obtains the same pressure model from \law{distributed}.

\paragraph{What the closed-form friction gives.} Independence of the triangulation and of the
decomposition, exactly and at no cost in mesh (Section~\ref{sec:ablation}). It removes a
requirement on the mesh and is not a gain in accuracy over a well-used quadrature. It is the
only element of \law{distributed} that we claim as new, and it is a direct extension of the
uniform-pressure result of Jia \cite{Jia1997} with the same primitives.
